\section{Analyse du problème - Baseline}

\subsection{État de l'art}


L'implémentation actuelle initialise d'abord l'environnement dans \texttt{fractal\_land.cpp} en construisant
 une carte fractale par raffinement récursif de sous-grilles. Dans \texttt{ant\_simu.cpp}, cette carte est 
 ensuite normalisée pour ramener les altitudes dans l'intervalle $[0,1]$, ce qui fixe le coût de déplacement
  utilisé par les fourmis pendant la simulation.


La population est ensuite instanciée avec un nombre fixe de fourmis, des positions initiales pseudo-aléatoires
 et des paramètres globaux du modèle (\textit{exploration}, \textit{bruit}, \textit{évaporation}). La structure
  \texttt{pheronome} (dans \texttt{pheronome.hpp}) maintient deux valeurs de phéromones par cellule, un tampon
   de mise à jour, ainsi que des bordures fantômes marquées à \(-1\) pour empêcher les traversées hors domaine.


À chaque itération, la boucle principale de \texttt{ant\_simu.cpp} traite les événements SDL puis exécute 
\texttt{advance\_time}. Cette routine appelle \texttt{ant::advance} pour chaque agent, applique 
l'évaporation globale des phéromones (\texttt{do\_evaporation}) puis valide l'état suivant de la carte
 (\texttt{update}). L'évolution individuelle d'une fourmi, définie dans \texttt{ant.cpp}, combine un 
 choix local de voisin (aléatoire ou guidé par le maximum de phéromone), l'accumulation d'un budget de 
 mouvement (\texttt{consumed\_time < 1}), le marquage local des phéromones et la mise à jour de l'état 
 chargé/non chargé.


Le rendu de l'état simulé est assuré par \texttt{renderer.cpp} et \texttt{window.cpp}. Le rendu affiche 
le terrain, la position des fourmis, la distribution des deux phéromones et la courbe de nourriture cumulée. 
La couche fenêtre/renderer SDL est construite avec mécanismes de repli (\textit{fallback}) pour conserver 
l'exécution même lorsque certaines options matérielles ne sont pas disponibles.



\subsection{Profilage et identification des goulots d'étranglement}

\noindent
Le profilage a pour objectif d'isoler les coûts dominants de la simulation avant toute
optimisation, en particulier avant les étapes de vectorisation et de parallélisation.
Nous presentons ici la définition des
blocs à mesurer, instrumentation temporelle, protocole d'exécution reproductible et
structure des résultats attendus. Cette section décrit donc la méthode de profilage et les
objets techniques mobilisés, sans présenter de résultats chiffrés.

\begin{table}[H]
\centering
\caption{Découpage des parties mesurées pour le profilage de la simulation}
\begin{tabularx}{\textwidth}{@{}l l X l@{}}
\toprule
\textbf{ID} & \textbf{Nature} & \textbf{Rôle} & \textbf{Périmètre} \\
\midrule
P0 & Initialisation & Génération fractale du terrain & One-shot (démarrage) \\
P1 & Initialisation & Normalisation des coûts du terrain & One-shot (démarrage) \\
P2 & Initialisation & Création et placement de la population & One-shot (démarrage) \\
K0 & Noyau simulation & Coût total du noyau par itération & Itératif \\
K1 & Noyau simulation & Avance globale de toutes les fourmis & Itératif \\
K2 & Noyau simulation & Logique interne hors marquage phéromones & Itératif \\
K3 & Noyau simulation & Marquage local des phéromones & Itératif \\
K4 & Noyau simulation & Évaporation globale des phéromones & Itératif \\
K5 & Noyau simulation & Validation/swap de la carte de phéromones & Itératif \\
R0 & Rendu SDL & Affichage (terrain, fourmis, phéromones, courbe) & Itératif, hors noyau \\
E0 & Événements SDL & Traitement de \texttt{SDL\_PollEvent} & Itératif, hors noyau \\
\bottomrule
\end{tabularx}
\end{table}

\noindent
Les relations structurelles utilisées pour l'analyse des goulots sont explicites:
\[
K0 = K1 + K4 + K5
\]
\[
K1 = K2 + K3 + \text{reste}(K1)
\]
Les blocs \(P0\), \(P1\) et \(P2\) sont des coûts \textit{one-shot} d'initialisation et ne doivent
pas être mélangés aux métriques par itération. Les blocs \(R0\) (rendu SDL) et \(E0\)
(gestion des événements SDL) sont mesurés séparément et traités comme non parallélisables
dans l'analyse du noyau de calcul.

\noindent
La stratégie instrumentale côté C++ s'appuie sur \texttt{timing\_profile.hpp} et
\texttt{timing\_profile.cpp}. Le profilage est désormais \emph{instance-based} (plus d'état
mutable global partagé): chaque exécution benchmark manipule son objet \texttt{TimingProfile}
et les compteurs d'itération \texttt{IterTimingNs}. \texttt{timing\_profile.hpp} fournit
l'horloge monotone (\texttt{profile\_now\_ns}) utilisée pour les frontières de kernels
(\(K0\), \(K1\), \(K4\), \(K5\)), tandis que \(K2/K3\) sont accumulés côté boucle agents
comme temps de travail (\emph{work-time}) dans les métriques d'itération.
Cette organisation sépare proprement coûts d'initialisation, coûts noyau et coûts hors noyau.

\noindent
L'automatisation est assurée par \texttt{scripts/measure\_temps.sh}. Le pipeline suit une
séquence reproductible: compilation release, puis 5 répétitions d'exécution avec 1200
itérations par répétition, 200 itérations de \textit{warm-up} exclues et 1000 itérations
effectivement mesurées, en configuration \texttt{OMP\_NUM\_THREADS=1}. Le script consolide les
artefacts méthodologiques (\texttt{raw\_measurements.csv}, \texttt{init\_times.csv},
\texttt{table\_a.csv}, \texttt{table\_b.csv}, \texttt{amdahl\_prediction.csv},
\texttt{summary.md}) afin de préparer la lecture des goulots. Les artefacts générés sont:
\begin{itemize}[leftmargin=*]
    \item \texttt{raw\_measurements.csv}: mesures brutes par répétition (compteurs en ns et valeurs dérivées par itération).
    \item \texttt{init\_times.csv}: synthèse des coûts \textit{one-shot} d'initialisation (\(P0\), \(P1\), \(P2\)).
    \item \texttt{table\_a.csv}: métriques des blocs de profilage (moyenne, dispersion et part relative du noyau).
    \item \texttt{table\_b.csv}: indicateurs globaux par répétition (\(T\_\text{iter}\), \(T\_\text{kernel}\), \(T\_\text{render}\), débit).
    \item \texttt{amdahl\_prediction.csv}: valeurs théoriques de speedup issues du modèle d'Amdahl à partir de la fraction parallélisable estimée.
    \item \texttt{summary.md}: synthèse lisible du protocole et des tableaux générés.
\end{itemize}

\noindent
La méthode de lecture priorise les métriques \texttt{ms/iter} et \% noyau pour identifier les
\textit{hotspots} dominants. Les conclusions doivent d'abord cibler les blocs du noyau
\((K0\text{--}K5)\), puis isoler l'impact des composants hors noyau \((R0,E0)\) pour éviter
les biais d'interprétation.

\subsection{Résultats}

\noindent
Les résultats présentés ci-dessous proviennent de la campagne de mesure:\\
\texttt{results/test\_time\_measurements\_local2/aos/latest/summary.md}. La configuration utilisée
est la suivante: 5 répétitions, 1200 itérations par répétition, 200 itérations de
\textit{warm-up}, 1000 itérations mesurées, et exécution avec \texttt{OMP\_NUM\_THREADS=1}.

\begin{table}[H]
\centering
\caption{Temps one-shot d'initialisation}
\label{tab:results-init}
\begin{tabular}{@{}lrr@{}}
\toprule
\textbf{Partie} & \textbf{ms moyen} & \textbf{écart-type} \\
\midrule
P0 & 8.222398 & 1.107389 \\
P1 & 0.830145 & 0.541537 \\
P2 & 0.171107 & 0.032904 \\
\bottomrule
\end{tabular}
\end{table}

\begin{table}[H]
\centering
\caption{Décomposition des coûts (profilage de base)}
\label{tab:results-decomp}
\begin{tabular}{@{}lrrr@{}}
\toprule
\textbf{Partie} & \textbf{ms/iter moyen} & \textbf{écart-type} & \textbf{\% kernel} \\
\midrule
K1 & 7.583137 & 0.803269 & 94.3004 \\
K2 & 3.278775 & 0.344471 & 40.7733 \\
K3 & 2.662628 & 0.413677 & 33.1112 \\
K4 & 0.447192 & 0.087824 & 5.5611 \\
K5 & 0.010333 & 0.002728 & 0.1285 \\
R0 & 2.190812 & 0.341754 & n/a \\
E0 & 0.005376 & 0.000923 & n/a \\
\bottomrule
\end{tabular}
\end{table}

\begin{table}[H]
\centering
\small
\caption{Indicateurs globaux par répétition}
\label{tab:results-global}
\begin{tabular}{@{}lrrrr@{}}
\toprule
\textbf{Répétition} & \textbf{T\_iter\_total (ms)} & \textbf{T\_kernel (ms)} & \textbf{T\_render (ms)} & \textbf{fourmis/s} \\
\midrule
1 & 8.619821 & 6.832405 & 1.783345 & 731806.773612 \\
2 & 11.077200 & 8.649231 & 2.422158 & 578086.088303 \\
3 & 9.838528 & 7.800593 & 2.032789 & 640976.879203 \\
4 & 9.514154 & 7.527849 & 1.981301 & 664200.401132 \\
5 & 12.138599 & 9.397280 & 2.734469 & 532068.839698 \\
Moyenne & 10.237660 & 8.041472 & 2.190812 & 629427.796390 \\
Écart-type & 1.234388 & 0.893450 & 0.341754 & 69193.508032 \\
\bottomrule
\end{tabular}
\end{table}

\noindent
Fraction parallélisable candidate:
\[
F = (K1 + K4) / K0 = 0.997839
\]

\begin{table}[H]
\centering
\caption{Prédiction théorique selon Amdahl}
\label{tab:results-amdahl}
\begin{tabular}{@{}rr@{}}
\toprule
\textbf{Threads} & \textbf{S(p)} \\
\midrule
2 & 1.995687 \\
4 & 3.974235 \\
8 & 7.880787 \\
16 & 15.497644 \\
\bottomrule
\end{tabular}
\end{table}

\subsubsection{Analyse des Résultats}

\noindent
Cette sous-section interprète les mesures présentées dans les tableaux~\ref{tab:results-init},
~\ref{tab:results-decomp},~\ref{tab:results-global} et~\ref{tab:results-amdahl}, ainsi que les
indicateurs \(F\) et \(S(p)\).
L'objectif est d'identifier les priorités de parallélisation à partir des données observées,
en séparant clairement ce qui relève des coûts d'initialisation, du noyau de simulation et des
bornes théoriques de gain.

\begin{enumerate}[leftmargin=*]
    \item \textbf{Coûts d'initialisation (\(P0\text{--}P2\)).}
    Dans le tableau~\ref{tab:results-init}, \(P0\) domine très largement la phase
    d'initialisation (\(8.222398\) ms), alors que \(P1\)
    (\(0.830145\) ms) et \(P2\) (\(0.171107\) ms) restent faibles. Cette hiérarchie est cohérente
    avec le fait que la génération fractale du terrain est la tâche la plus lourde au démarrage.
    En revanche, comme \(P0\text{--}P2\) sont des coûts \textit{one-shot} hors boucle itérative,
    ils ne constituent pas la priorité principale pour améliorer le temps de simulation par itération.

    \item \textbf{Priorité de parallélisation au niveau noyau.}
    D'après le tableau~\ref{tab:results-decomp}, \(K1\) représente \(94.3004\%\) du temps noyau,
    ce qui en fait la cible prioritaire.
    Le raisonnement est direct: réduire \(K1\) a l'impact potentiel maximal sur \(K0\),
    avant \(K4\) (évaporation) et \(K5\) (mise à jour de la carte de phéromones), dont les
    contributions relatives sont nettement plus faibles.

    \item \textbf{Décomposition interne de \(K1\): \(K2\) avant \(K3\).}
    Toujours dans le tableau~\ref{tab:results-decomp}, \(K2\) pèse \(40.7733\%\) du noyau contre
    \(33.1112\%\) pour \(K3\).
    L'ordre de priorité est donc de paralléliser d'abord \(K2\), puis \(K3\).
    Cela confirme que la logique de déplacement/décision des fourmis est plus coûteuse que le
    marquage des phéromones.

    \item \textbf{Stabilité expérimentale inter-répétitions.}
    Le tableau~\ref{tab:results-global} montre une faible dispersion entre les 5 répétitions
    (\(T\_\text{iter\_total}\), \(T\_\text{kernel}\), \(T\_\text{render}\), \textit{fourmis/s}),
    ce qui indique une bonne stabilité de la campagne de mesure. Cette stabilité provient du protocole
    (phase de \textit{warm-up} exclue et répétitions multiples), et non d'un effet spécifique
    ``à partir de la cinquième itération''.

    \item \textbf{Interprétation de la fraction parallélisable \(F\).}
    La valeur \(F = 0.998797\), proche de 1, n'est pas anormale dans ce cadre:
    la définition \(F=(K1+K4)/K0\) inclut les blocs dominants du noyau et exclut les
    parties hors noyau (\(R0\), \(E0\)). Cette valeur doit toutefois rester interprétée
    comme une fraction candidate théorique pour le noyau, et non comme une garantie de
    speedup réel, car les surcoûts de synchronisation, d'accès mémoire et d'équilibrage
    peuvent réduire le gain observé.

    \item \textbf{Lecture correcte de la prédiction d'Amdahl.}
    Le tableau~\ref{tab:results-amdahl} montre un potentiel de gain important avec l'augmentation
    du nombre
    de threads. En particulier, \(S(16)=15.497644\) correspond à un facteur multiplicatif
    d'environ \(\times 15.7\), et non à \(15\%\). Cette tendance soutient clairement la
    stratégie de parallélisation; la suite du projet doit donc viser d'abord \(K1\), puis \(K4\),
    pour se rapprocher de ce plafond théorique.
\end{enumerate}

\noindent
En synthèse, les deux \textit{hotspots} centraux sont \(K1\) au niveau global et \(K2\) à
l'intérieur de \(K1\). Les mesures actuelles établissent une base robuste pour prioriser les
optimisations, mais l'écart entre potentiel théorique et performance finale devra être vérifié
expérimentalement après implémentation de la parallélisation. La feuille de route est donc
cohérente: concentrer l'effort d'abord sur les blocs dominants du noyau afin de maximiser le
retour de chaque évolution de code.
